\documentclass[a4,11pt]{aleph-notas}

% -- Paquetes adicionales
\usepackage{enumitem}
\usepackage{aleph-comandos}

% -- Datos
\institucion{Facultad de Ciencias Exactas, Naturales y Ambientales}
\carrera{Catálogo STEM}
\asignatura{Métodos Numéricos}
\tema{Clase 00: Introducción a los Métodos Numéricos}
\autor{Mario Cueva - Andrés Merino}
\fecha{Periodo 2026-1}

\logouno[0.18\textwidth]{Logos/logoPUCE_04_ac}
\definecolor{colortext}{HTML}{1957d1}
\definecolor{colordef}{HTML}{1957d1}
\fuente{montserrat}


\begin{document}

\encabezado

%%%%%%%%%%%%%%%%%%%%%%%%%%%%%%%%%%%%%%%%
\section*{Resultado de Aprendizaje}
%%%%%%%%%%%%%%%%%%%%%%%%%%%%%%%%%%%%%%%%

%%%%%%%%%%%%%%%%%%%%%%%%%%%%%%%%%%%%%%%%
\subsection*{RdA de la asignatura:}
\begin{itemize}[leftmargin=*]
    \item \textbf{RdA 1:} Comprender los conceptos fundamentales de los Métodos Numéricos, incluyendo la solución de ecuaciones y sistemas, diferenciación e integración numérica y la optimización numérica, destacando su importancia en el análisis y resolución de problemas matemáticos.
    \item \textbf{RdA 2:} Computar cálculos relacionadas con los Métodos Numéricos, incluyendo la solución de ecuaciones y sistemas, la diferenciación e integración numéricas, y la optimización numérica, para generar aproximaciones a soluciones.
    \item \textbf{RdA 3:} Aplicar los conceptos y técnicas de los Métodos Numéricos en la resolución de problemas reales, incluyendo la solución de ecuaciones y sistemas, la diferenciación e integración numéricas, y la optimización numérica, demostrando habilidad para abordar desafíos prácticos y situaciones del mundo real.
\end{itemize}

%%%%%%%%%%%%%%%%%%%%%%%%%%%%%%%%%%%%%%%%
\subsection*{RdA de la actividad:}
\begin{itemize}[leftmargin=*]
    \item Reconocer el propósito de los métodos numéricos como herramientas para obtener aproximaciones útiles cuando no se dispone de una solución exacta simple.
    \item Identificar los principales bloques temáticos del curso y su relación con problemas aplicados.
    \item Establecer acuerdos de trabajo para el uso de recursos, actividades, evaluaciones y herramientas computacionales durante el periodo.
\end{itemize}

%%%%%%%%%%%%%%%%%%%%%%%%%%%%%%%%%%%%%%%%
\section*{Introducción}
%%%%%%%%%%%%%%%%%%%%%%%%%%%%%%%%%%%%%%%%

%%%%%%%%%%%%%%%%%%%%%%%%%%%%%%%%%%%%%%%%
\paragraph{Pregunta inicial:} 
¿Qué haríamos si un problema importante no se puede resolver con una fórmula exacta, pero necesitamos una respuesta confiable hoy?

\paragraph{Ejemplos de partida:}
\begin{itemize}[leftmargin=*]
    \item Estimar la raíz cuadrada de $2$ cuando solo se necesitan cuatro decimales.
    \item Calcular el área bajo una curva cuando la antiderivada no es sencilla.
    \item Predecir la temperatura de un día usando datos medidos en horas específicas.
    \item Encontrar el precio que maximiza la ganancia de un producto con restricciones reales.
\end{itemize}

%%%%%%%%%%%%%%%%%%%%%%%%%%%%%%%%%%%%%%%%
\section*{Desarrollo}
%%%%%%%%%%%%%%%%%%%%%%%%%%%%%%%%%%%%%%%%

%%%%%%%%%%%%%%%%%%%%%%%%%%%%%%%%%%%%%%%%
\subsection*{Actividad 1: ¿Para qué sirven los métodos numéricos?}

La actividad presenta el sentido del curso mediante una conversación guiada y ejemplos aplicados.

\paragraph{¿Cómo lo haremos?}  
\begin{itemize}[leftmargin=*]
    \item \textbf{Diagnóstico inicial:} cada estudiante escribe un problema de su carrera donde una respuesta aproximada pueda ser útil.
    \item \textbf{Discusión guiada:} se clasifican los problemas propuestos según involucren ecuaciones, datos, cambios, acumulación u optimización.
    \item \textbf{Clase interactiva:} se explica la diferencia entre solución exacta, aproximación, error y tolerancia.
    \item \textbf{Síntesis:} se construye una definición breve de método numérico a partir de las respuestas del grupo.
\end{itemize}

\paragraph{Ejemplo guiado:}
Se desea estimar $\sqrt{2}$ para diseñar una pieza donde basta una precisión de dos decimales.
\[
    \sqrt{2}=1.41421356\ldots \approx 1.41.
\]
La aproximación $1.41$ no es exacta, pero puede ser adecuada si la tolerancia del problema permite errores menores que $0.01$.

\paragraph{Actividad corta:}
\begin{enumerate}[leftmargin=*]
    \item Escriba una cantidad que haya medido o estimado recientemente.
    \item Indique si esa cantidad era exacta o aproximada.
    \item Proponga una tolerancia razonable para esa cantidad.
\end{enumerate}

%%%%%%%%%%%%%%%%%%%%%%%%%%%%%%%%%%%%%%%%
\subsection*{Actividad 2: Aproximación, error y tolerancia}

La actividad muestra que una aproximación debe juzgarse según el tamaño del error y el contexto del problema.

\paragraph{¿Cómo lo haremos?}
\begin{itemize}[leftmargin=*]
    \item \textbf{Ejemplo numérico:} se compara el valor verdadero $x=3.1416$ con la aproximación $\widehat{x}=3.14$.
    \item \textbf{Cálculo guiado:} se calcula el error absoluto $E_a=|x-\widehat{x}|=0.0016$.
    \item \textbf{Interpretación:} se analiza si el error es aceptable para un cálculo escolar, una medición de laboratorio o una aplicación médica.
    \item \textbf{Trabajo individual:} cada estudiante resuelve un cálculo de error con una aproximación distinta.
\end{itemize}

\paragraph{Ejemplos para resolver en clase:}
\begin{enumerate}[leftmargin=*]
    \item Si el valor verdadero es $x=10$ y la aproximación es $\widehat{x}=9.8$, calcule el error absoluto.
    \item Si el valor verdadero es $x=0.05$ y la aproximación es $\widehat{x}=0.04$, calcule el error absoluto y discuta si es pequeño.
    \item Si dos aproximaciones consecutivas son $x_1=1.41$ y $x_2=1.414$, calcule $|x_2-x_1|$.
\end{enumerate}

\paragraph{Pregunta para discusión:}
¿Puede un mismo error absoluto ser aceptable en un problema y peligroso en otro?

%%%%%%%%%%%%%%%%%%%%%%%%%%%%%%%%%%%%%%%%
\subsection*{Actividad 3: Mapa del curso mediante mini-casos}

La actividad conecta cada unidad del curso con un problema aplicado que se resolverá con aproximaciones.

\paragraph{¿Cómo lo haremos?}
\begin{itemize}[leftmargin=*]
    \item \textbf{Representación de números:} se analiza por qué una computadora no siempre representa exactamente los decimales.
    \item \textbf{Resolución de ecuaciones:} se plantea encontrar la temperatura donde dos modelos físicos coinciden.
    \item \textbf{Sistemas de ecuaciones:} se modela una mezcla de soluciones con concentraciones diferentes.
    \item \textbf{Interpolación:} se estima un dato faltante usando mediciones cercanas.
    \item \textbf{Diferenciación e integración numérica:} se aproxima velocidad, acumulación o área a partir de datos discretos.
    \item \textbf{Optimización:} se busca el valor que maximiza una ganancia o minimiza un costo.
\end{itemize}

\paragraph{Actividad en grupos:}
Cada grupo elige uno de los mini-casos, identifica la pregunta matemática principal y propone qué tipo de resultado aproximado esperaría obtener.

%%%%%%%%%%%%%%%%%%%%%%%%%%%%%%%%%%%%%%%%
\subsection*{Actividad 4: Primer contacto computacional}

La actividad introduce el uso de calculadora o computador para observar el efecto del redondeo.

\paragraph{¿Cómo lo haremos?}
\begin{itemize}[leftmargin=*]
    \item \textbf{Exploración:} se calcula $1/3$, $2/3$ y $1/3+2/3$ en una calculadora o computador.
    \item \textbf{Observación:} se compara el resultado mostrado con el valor exacto esperado.
    \item \textbf{Reflexión:} se discute por qué la cantidad de dígitos disponibles afecta la respuesta.
    \item \textbf{Registro:} cada estudiante anota una conclusión sobre precisión y redondeo.
\end{itemize}

\paragraph{Verificación de aprendizaje:}
\begin{enumerate}[leftmargin=*]
    \item Explique con una oración qué estudian los métodos numéricos.
    \item Mencione dos temas del curso que le parezcan útiles para su carrera.
    \item Describa una situación real donde una aproximación sea más práctica que una solución exacta.
    \item Explique por qué una tolerancia debe elegirse antes de aceptar una aproximación.
\end{enumerate}


%%%%%%%%%%%%%%%%%%%%%%%%%%%%%%%%%%%%%%%%
\section*{Cierre}
%%%%%%%%%%%%%%%%%%%%%%%%%%%%%%%%%%%%%%%%

\paragraph{Tarea:}
Leer el \href{https://fcena-puce.github.io/FCENA-PUCE/MetodosNumericos-05-N0062/2-Resumenes/Resumen01.pdf}{Resumen01.pdf} e identificar un ejemplo cotidiano donde aparezcan errores de redondeo o medición.

\paragraph{Pregunta de investigación:}  
\begin{enumerate}[leftmargin=*]
    \item ¿Qué áreas profesionales dependen de aproximaciones numéricas para tomar decisiones?
    \item ¿Qué riesgos pueden aparecer cuando se ignora el error en un cálculo numérico?
    \item ¿Cómo ayuda la programación a implementar métodos numéricos de forma eficiente?
\end{enumerate}
    
\paragraph{Para la próxima clase:}  
Traer una calculadora o computador y repasar operaciones básicas con potencias, notación científica y porcentajes.

\end{document} 
